\chapter{Maintenance Comments}\label{sec:MaintenanceComments}
It is very likely that source code is modified more than once during its lifetime: bugs are fixed, functionality is added or removed, refactorings are applied, or the code is migrated to other platforms or newer versions of the programming language, of the underlying libraries, of the connected systems and/or of the operating system.


\bdq{\nameref{sec:TagSince}}


There is a practice that all the changes made in the code will be commented in the code. These comments are usually referred to as “Maintenance Comments”.

This may look like this:\footnote{The abbreviation “FSP” stands for “Field Service Patch” …}
\begin{lstlisting}
…
//<<BEGIN FSP-0815 – applied by Mirabel Madrigal
<Additional Code>
…
//>>END FSP-0815
…
//<<BEGIN FSP-4711 – applied by Micky Mouse
//<Old Code>
//…
//>><<
<New Code>
…
//>>END FSP-4711
…
//<<BEGIN FSP-4917 – applied by Donald Duck
//<Old Code>
//…
//>>
<Unmodified Code>
…
//<<
<New Code>
…
//>>END FSP-4917
…
\end{lstlisting}

Also it looks good at a first glance, it has proved that this is not a good practice at all, and a really bad practice if dealing with code that is managed by an SCCS.

First, this practice causes problems for the compare tools coming with the SCCS – at least it will make it more difficult to read the comparison results from those tools.\footnote{Most comparison tools will recognise the out-commenting of the old code as a change and the new code as additional code instead of the replacement for the old code. This is at least confusing when a code revisions are made on the fix.} And the main function of those comments – documenting the changes – is much better served by the SCCS tools themselves.

Second, this gets even worse when the fix needs a fix. Just think about overlapping changes, like FSP-4712 changes lines from the new code of FSP-4711, together with lines directly below, or changes that has to be partially reverted, and so on.

And finally, the comments does not help to identify if your are currently running the patched code, or the old one.

On the other hand, nothing can be said against adding a line to the file or class comment that lists the patches that were applied to the code.

Therefore I recommend to introduce an annotation that can be used to mark elements that are affected by a fix.

Such an annotation can provide the BUG number of the fix, together with a short description of the issue. This annotation replaces a comment about the applied patches, with the advantage that it can be retrieved also from the compiled classes, even at runtime.

Chapter \tqfullvref{sec:PatchIdentification} provides an example of such an annotation.

Old code will be removed and not commented out, or just replaced by the new code. One or the other short comment with a hint is not mandatory, but does not harm either.

When using the suggested annotation from chapter \tqref{sec:PatchIdentification}, this could look like this:
\begin{lstlisting}
…
@BUG( id = "BUG-0815", comment = "Superseded by BUG-4711" )
@BUG( id = "BUG-4711", comment = "No end criterion for loop" )
@BUG( id = "BUG-4712", comment = "Exception was swallowed" )
public final void myMethod()
{
    ForeverLoop: while( true )
    {
        try
        {
            …
            
            if( !hasMore() ) break ForeverLoop; // BUG-4711
        }
        catch( final IllegalStateException e )
        {
            log( e );
            break ForeverLoop; // BUG-4712: Added exception handling
        }
    }   //  ForeverLoop:
}   //  myMethod()
…
\end{lstlisting}

The relevant details about the fix/change/enhancement can now be found in the SCCS and/or in your Bug Tracking system, referenced through the id of the annotation.

%==============================================================================
% End of file!!
%==============================================================================

